\documentclass{article}
\usepackage[utf8]{inputenc}
\usepackage{amsmath, amssymb}
\usepackage{array}

\begin{document}
	
	\section*{Problema 2: Pensamiento Matemático - Psicología}
	
	\subsection*{a) Proposiciones simples}
	
	Definimos las proposiciones:
	
	
	\[
	\begin{aligned}
		p &: \text{La persona sigue las estrategias terapéuticas indicadas.} \\
		\sim p &: \text{La persona no sigue las estrategias terapéuticas indicadas.} \\
		q &: \text{La persona presenta dificultades emocionales.} \\
		r &: \text{La persona logra mejorar su bienestar psicológico.}
		\sim r &: \text{La persona no logra mejorar su bienestar psicológico.}
	\end{aligned}
	\]
	
	
	
	El enunciado:  
	\textit{“Si una persona no sigue las estrategias terapéuticas indicadas, entonces presenta dificultades emocionales o no logra mejorar su bienestar psicológico.”}
	
	En lenguaje simbólico:
	
	
	\[
	\sim p \to (q \, \lor \sim r)
	\]
	
	
	
	\subsection*{b) Tabla de verdad}
	
	La proposición compuesta es:
	
	
	\[
	\sim p \to (q \, \lor \sim r)
	\]
	
	
	
	\begin{center}
		\begin{tabular}{|c|c|c|c|c|c|c|}
			\hline
			$p$ & $q$ & $r$ & $\sim p$ & $\sim r$ & $q \lor \sim r$ & $\sim p \to (q \, \lor  \sim r)$ \\
			\hline
			V & V & V & F & F & V \\
			V & V & F & F & V & V \\
			V & F & V & F & F & V \\
			V & F & F & F & V & V \\
			F & V & V & V & F & V \\
			F & V & F & V & V & V \\
			F & F & V & V & F & F \\
			F & F & F & V & V & V \\
			\hline
		\end{tabular}
	\end{center}
	
	\subsection*{Conclusión}
	
	La proposición es verdadera en la mayoría de los casos, excepto cuando la persona no sigue las estrategias terapéuticas (\(\sim p = V\)) y no presenta dificultades emocionales ni falta de bienestar (\(q \lor \sim r = V\)).  
	Por lo tanto, la proposición es una:
	
	
	\[
	\boxed{\text{Contingencia}}
	\]
	
	
	
\end{document}
